% Generated draft: requires mathematical and source review.
\documentclass[11pt]{article}
\usepackage[margin=25mm]{geometry}
\usepackage{fontspec}
\setmainfont{Latin Modern Roman}
\usepackage{amsmath,amssymb,mathtools,unicode-math}
\setmathfont{Latin Modern Math}
\usepackage{enumitem,needspace,xurl,hyperref}
\setlength{\parindent}{0pt}
\setlength{\parskip}{5pt}
\setlength{\emergencystretch}{3em}
\newcommand{\N}{\mathbb{N}}
\newcommand{\Z}{\mathbb{Z}}
\newcommand{\Q}{\mathbb{Q}}
\newcommand{\R}{\mathbb{R}}
\newcommand{\C}{\mathbb{C}}
\newcommand{\F}{\mathbb{F}}
\newcommand{\A}{\mathbb{A}}
\newcommand{\PP}{\mathbb{P}}
\newcommand{\E}{\mathbb{E}}
\newcommand{\eps}{\varepsilon}
\newcommand{\dd}{\,\mathrm{d}}
\begin{document}
\section*{UnsolvedMath ID: 6700016}
\textbf{Source problem code:} catalogue entry
\subsection{Definitions and mathematical statement}
The conjecture concerns hypersurfaces in an ambient \(n\)-manifold, hence submanifolds or currents of dimension \(n-1\), not arbitrary codimension. A natural homological formulation fixes a smooth closed \(X^n\) and considers the space \(\operatorname{Met}^\infty(X)\) of smooth Riemannian metrics with its \(C^\infty\) topology. For integral cycles, an area-minimizer in a class \(a\in H_{n-1}(X;\Z)\) is a current \(T\) with \([T]=a\) and
\[
\mathbf M_g(T)\le\mathbf M_g(S)\qquad\text{for every integral cycle }S\text{ with }[S]=a.
\]
Nonsingularity means that its support has no singular points; multiplicity one is not silently required.

In this formulation, define
\[
\mathcal R(X)=\{g\in\operatorname{Met}^\infty(X):
 \text{every such homological minimizer has empty singular set}\}.
\]
A generic-metric version asks that \(\mathcal R(X)\) contain a residual subset, meaning a countable intersection of open dense sets in the metric space of metrics. The weaker arbitrarily-small-perturbation statement asks only
\[
\forall g\ \forall\text{ neighborhood }\mathcal U\ni g,
\qquad \mathcal U\cap\mathcal R(X)\ne\varnothing.
\]
Residual genericity and density are not identical assertions. The extracted source uses both expressions without specifying topology or exact quantifiers. Nor does it fix closedness, integral versus mod-two coefficients, homological minimizers versus fixed-boundary or local minimizers, or whether one class or all classes are treated simultaneously. The displayed homological model is therefore explicit but conditional on those scope choices.

The source attributes the belief to Brian White and reports confirmation by Nathan Smale in 1993 for ambient dimension \(n=8\), with a stated extension of the Schoen--Yau theorem to that dimension. The research summary also mentions later Chodosh--Li generic-regularity progress, without a precise theorem; it reports the all-dimensional statement unresolved. Neither its loose all-codimension wording nor its unspecified later results are promoted beyond the hypersurface question.
\subsection{Short English statement}
Should generic smooth metrics make volume-minimizing hypersurfaces nonsingular, as reported in ambient dimension eight? The minimization class, coefficient group, topology, and density versus residual-generic quantifiers need to be fixed.
\subsection{Sources}
\begin{itemize}
\item Local UnsolvedMath registry, permanent ID \(6700016\), exact source code \texttt{catalogue entry}, scalar-curvature item [?16], ``Singularities are Unstable.'' Proposer Misha Gromov, 2017; conjectural belief attributed in the source to Brian White.
\item Gromov, \emph{101 Questions, Problems and Conjectures around Scalar Curvature} (2017), printed tag [?16], occurrence 16, recorded PDF page 21; \url{https://www.ihes.fr/~gromov/wp-content/uploads/2018/08/101-problemsOct1-2017.pdf}. Recorded access 2026-07-29; pdf-text-normalized extraction.
\item Nathan Smale, ``Generic regularity of homologically area minimizing hypersurfaces in eight-dimensional manifolds,'' \emph{Communications in Analysis and Geometry} (1993), as supplied in the research summary; source citation [118]. Complete volume/page details are absent.
\item Supplied PARTIAL-PROGRESS research summary additionally mentions Chodosh--Li (2021) without a title, precise statement, or hypotheses. That attribution requires verification before assigning a higher-dimensional conclusion.
\end{itemize}
\end{document}
